% !TeX root = ../papers/nc-p3.tex
\section{片側対称枝の正則性を測る二つの交代形式}
\label{sec:single-edge-instanton-extension}
\begin{theorem}\label{thm:single-edge-instanton-extension}
四辺有限、片側対称、反対階数 $\rho>6$ の標準有限Frobenius窓を $U,V,U,Z$ の順に書く。
それぞれ $U,V$ と $U,Z$ の辺で、定理\ref{thm:finite-square-tangent-instantons}の
$U$ の正規化平方根の接空間を取る。各辺で次は同値である。
\begin{enumerate}
\item その辺のインスタントン層が局所自由である。
\item 接空間が階数10の交代形式を含む。
\end{enumerate}
成立するとき接空間は一次元で、全ての非零な方向が階数10である。
特に $\rho=10$ の枝の全次数正則性は二辺の交代形式の線型系がともに非退化形式を含むことに同値である。
有限点配置の任意の特殊値に適用できる。
\end{theorem}
\begin{proof}
一方の辺で $K=R^\perp$、モナドを $C=(L\otimes\mathcal O(-1)\xrightarrow{a}K\otimes\mathcal O
\xrightarrow{b}V^*\otimes\mathcal O(1))$ とする。
$L$ は $b$ の一次シジジー全体で次元4である（定理\ref{thm:finite-symmetric-instanton-monad}）。
接方向は交代写像 $T:K^*\to K$ で $b T b^*=0$ を満たすものに等しい。
実際、各 $u$ で $uV$ に制限した交代形式が零という条件が、偏極により $U$ で交代、
$V$ で対称という条件になる。従って $T b^*=aP$ を満たす定数写像 $P:V\to L$ が一意に存在し、
転置によってもう一つの端でも可換となる。これは $C^\vee\to C$ の鎖写像である
（双対複体の微分の符号を片端に入れる）。$T$ が可逆なら
$b^*$ は各点で単射だから $P$ は可逆、$a$ は部分束であり、層は局所自由である。

逆に束 $F$ であれば両複体は $F^*,F$ のモナドである。
これらの間の導来Homの次数零は大域鎖写像そのものに等しい。
実際内部Hom複体の直線束の捻りは $-2$ から2に入り高次コホモロジーがなく、
次数 $-1$ のホモトピーの大域切断もない。これらは同じモナド同士、双対モナド同士でも成立する。
$F$ は階数2、安定で行列式自明なので、$F^*\to F$ のHomは外積の逆同型の張る直線である。
この同型と逆同型は一意な鎖写像に持ち上がり、その合成は鎖写像として恒等となる。
交代性も一意性により持ち上がる。中間項の写像が可逆で交代となることは、
複体を次数 $-1,0,1$ に置くと次数0の項の転置にKoszul符号がかからないことからも分かる。
こうして非退化な $T$ が得られる。他の接方向も鎖写像を与え、中間項でその方向が読めるので
この直線に入る。最後は定理\ref{thm:finite-rank-ten-full-regularity}の必要十分条件による。
\end{proof}

\paragraph{第一辺からの有理的な正則延長。}
前節の明示行列の第一辺 $R$ は $R^{\rm op}$ を次の辺に取ると実際に全次数正則な四窓に延びる。
反転三次交叉は次元4で、その基底を最後の文字で束ねたものを $w\in U\otimes V\otimes U\otimes Z$
とすれば選択は $Z$ の基底変更だけである。$w$ は二つの $U$ で対称で、二つの隣接切断階数は6、
反対切断階数は10である。$U,Z$ の辺を $w$ の隣接像として構成する。
\path{computations/verify_single_edge_instanton_extension.py} は二つの辺の両射影の有限性、
交代形式の各解空間が非退化な一本であることを有理数上で確かめる。さらに全始点で、
部分収縮の次元 $4,6,4,1$ と、左右から一文字伸ばした空間同士の交叉次元 $4,1,0$ を
良い素数で調べる。収縮による包含と有理数上の切断階数が下界を与えるので、
これらの交叉の上界は標数零にも戻る。従って標準有限Frobenius条件が全て成立し、
上の定理で正則性が従う。$Z$ を $V$ と同一視して反転する候補は前節の逆向き三次交叉により除かれる。
